% =============== 导言区 ===============
\documentclass{ctexart}

% 设置页边距的宏包
\usepackage{geometry}
\geometry{ % 设置整体页面格式
	a4paper,
	left = 2.1cm,
	right = 2.1cm,
	bottom = 2.1cm,
	top = 2.5cm
}

% 设置章节标题格式
%\CTEXsetup[
%	name = {\S\ ,},
%	aftername = {\hspace{2ex}}
%	]{section}

% 目录和超链接
\usepackage{tocbibind}			% 添加参考文献到目录
\usepackage{hyperref}			% 构造引用超链接
\hypersetup{colorlinks = true}	% 超链接用彩色显示但不要边框

% AMS数学包
\usepackage{amsmath}
\usepackage{amsfonts}
\usepackage{amssymb}
\numberwithin{equation}{section}			% 设置公式编号包含节编号

% XeLaTeX logo
\usepackage{metalogo}	

% 使用颜色并添加预定义名和表格支持
\usepackage[dvipsnames,table]{xcolor}

% 支持语法高亮的代码环境：lstlisting
\usepackage{listings}
\lstset{ % 代码环境整体设置
	basicstyle = \ttfamily,
	keywordstyle = \bfseries,
	commentstyle = \rmfamily\upshape,
	stringstyle = \ttfamily\slshape,
	tabsize = 4,
	backgroundcolor = \color{lightgray},	% 需先引入颜色宏包
	frame = single,
	language = TeX							% 设置默认语言
}

% 优化脚注
\usepackage[perpage]{footmisc}	% 脚注每页清零
\usepackage{pifont}				% 优化带圈的数字
\renewcommand{\thefootnote}{\ding{\numexpr171+\value{footnote}}}

% 图表宏包
\usepackage{graphicx}			% 插图宏包\includegraphics
\usepackage{tabularx}			% 定宽表格
\usepackage{multirow}			% 合并行
\usepackage{booktabs}			% 三线表
\usepackage{float}				% 图表浮动体控制

% 参考文献宏包
\usepackage[super,square]{natbib}

% 定义定理类
\usepackage[thmmarks]{ntheorem}
\newtheorem{thm}{定理}[section]
{	% 利用分组使设置只对该分组内的定理类有效
	\theoremstyle{nonumberplain}
	\theoremheaderfont{\bfseries}
	\theorembodyfont{\normalfont}
	\theoremsymbol{\ensuremath{\Box}}
	\newtheorem{proof}{证明}
}

% 定义四元数乘法的简写
\newcommand{\quadprod}[2]
	{\ensuremath{\mathfrak{#1} \otimes \mathfrak{#2}}}

% 设置参考文献按引用顺序排序
\bibliographystyle{unsrtnat}

% 采用titlepage自定义页面，以下信息可注释
%\title{\LaTeX 基础}
%\author{XiaoCY}
%\date{\today}

% =============== 正文区 ===============
\begin{document}
%	\maketitle	% 由titlepage代替：自定义封面
	\begin{titlepage}
		\zihao{4}	% 设置封面采用四号字
		\ % 产生空格占位
		
		\vspace{\fill}
		
		\begin{flushright}
			{\Huge\bfseries\LaTeX 使用简介} \\
			\rule[1em]{\linewidth}{0.5ex}
			
			\begin{tabular}{rr}
				作者 & XiaoCY \\
				版本 & 1.1 \\
				完成日期 & 2020-02-06 \\
				最后修改 & 2020-02-07 \\
			\end{tabular}
		\end{flushright}
		
		\vspace{\stretch{3}}
		
		{\noindent\slshape Email: 
		 \href{mailto:chunyu2018@foxmail.com}
			{chunyu2018@foxmail.com}
		}
	\end{titlepage}

	\zihao{-4}				% 设置全文采用小四号字
	
	\pagenumbering{Roman}
	\tableofcontents		% 生成目录
	\clearpage
	
	\listoftables
	\clearpage
	
	\pagenumbering{arabic}
	\setcounter{page}{1}	% 正文重新开始于1
	
	\section{写在前面}
	\LaTeX 是一种区别于Word的排版软件，
	其格式由各种命令、环境控制，
	很容易做到内容与格式分离，
	这是它区别于Word的重要一点。
	这篇文章主要是本人学习 \LaTeX 的记录\cite{liuhaiyang}，
	收集常用的基本操作以便后期查阅，
	同时也为感兴趣的小伙伴们做一个粗略的介绍。
	
	本人学习时安装的发行版为 \TeX~Live，
	对于不知道如何入门的小伙伴，
	我同样推荐这个发行版，
	它避免了很多繁琐的配置。
	编译器以\XeLaTeX 为主，
	编辑器可采用TeXstudio。
	
	\section{基础介绍}
	\subsection{文档组成}
	\LaTeX 文档的格式通常是以后缀.tex结尾的文本文件，
	为了正常使用中文和Unicode的特殊符号，
	务必将文本以UTF-8编码进行保存。
	
	\LaTeX 文档可分为导言区和正文区。
	导言区在命令\verb|\begin{document}|之前，
	通常对文档的性质做一些设置，
	也可以自定义一些命令。
	导言区之后为正文区，
	是文档的主要内容。
	
	在导言区内，我们首先需要声明文档类，
	本文采用了命令\verb|\documentclass{ctexart}|来声明C\TeX 文档类。
	此外，很多优秀的宏包可以辅助我们对文档格式进行控制，
	为了使用这些宏包，
	我们需要在导言区使用命令\verb|\usepackage{<宏包名>}| 。
	（如无特殊说明，本文的命令示意中以尖括号连同内容表示命令，
	使用时不需要添加尖括号。）
	
	在正文区内，我们只需要输入文档的正文即可。
	编辑器内部的各种换行不会引起排版后文档的换行，
	中 文 文 档 的 多 余 空 格 也 不 会 在 文 档 中 出 现。 \par
	文字换行可用双斜线\verb|\\|实现，
	而换段则采用空行的形式，或者使用命令\verb|\par|。
	tex文件内的多个空行或多个\verb|\par|不会在正文引起多余的空行。
	换行和换段的区别是：
	当设置了段落首行缩进时，换行不会引入缩进。
	
	\subsection{使用符号}
	正文中绝大多数符号都可以由键盘直接输入，如@。
	但是有一些符号在 \LaTeX 中具有特殊的作用，
	因而不能直接在正文中使用。
	通常情况下我们可以在这些符号前面加上反斜线\textbackslash，
	但是也有个别例外。
	这些特殊符号的作用和使用方法如表 \ref{tab:symbol} 所示。
	
	\begin{table}
		\centering
		\caption{特殊符号及其说明}
		\label{tab:symbol}
		\begin{tabular}{cll}
			\toprule
			符号 & 作用 & 输入方式\\
			\midrule
			\~{} & 不可打断的空格 & \verb|\~{}| \\
			\# & 用于宏定义 & \verb|\#| \\
			\$ & 数学模式 & \verb|\$| \\
			\% & 注释符 & \verb|\%| \\
			\^{} & 上标 & \verb|\^{}| \\
			\& & 用于表格对齐 & \verb|\&| \\
			\{\ \} & 用于分组 & \verb|\{ \}| \\
			\_ & 数学下标 & \verb|\_| \\
			\textbackslash & 宏命令和转义符 & \verb|\textbackslash| \\
			\bottomrule
		\end{tabular}
	\end{table}
	
	除此之外，符号-在 \LaTeX 正文中有很多用途：
	在数学模式下它是减号，如$3-1=2$；
	单独使用时它是英文连字符，如good-looking；
	两个连用（-\,-）时用来表示数字范围，如1--10；
	三个连用（-\,-\,-）时是破折号---比如这里。
	
	为了在正文中使用空格，可以采用反斜线和空格组合实现。
	此外，有一种称为幻影的神奇空格，
	它由命令\verb|\phantom{<内容>}|生成，
	空格的长度与\verb|<内容>|所占据的长度相同。
	
	\subsection{设置字体}
	字体具有五种不同的性质，
	在 \LaTeX 中一起决定了文字最终的输出效果。
	字号（font size）是指文字的大小，
	常常被独立出来，看作不同于字体的单独的性质；
	字体编码（font encoding）指字体包含的符号，
	一般情况下不直接进行设定。
	使用最多的是其他的三个性质：
	字体族（font family）、字体形状（font shape）、
	字体系列（font series）\footnote{通常指字体的粗细和宽度}。
	字体的这三个性质也常称为字体的坐标。
	
	\LaTeX 提供了各种命令来对字体进行修改。
	命令和效果如表 \ref{tab:fontcoord} 所示。
	
	\begin{table}[hbp]
		\centering
		\caption{字体的坐标}
		\label{tab:fontcoord}
		\begin{tabular}{llll}
			\toprule
			字体族 & 带参数的命令 & 申明命令 & 效果 \\
			罗马 & \verb|\textrm{<text>}| & \verb|\rmfamily|
			    & \rmfamily Roman font family \\
			无衬线 & \verb|\textsf{<text>}| & \verb|\sffamily|
			    & \sffamily Sans serif font family \\
			打字机 & \verb|\texttt{<text>}| & \verb|\ttfamily|
			    & \ttfamily Typewriter font family \\
			\midrule
			字体形状 & 带参数的命令 & 申明命令 & 效果 \\
			直立 & \verb|\textup{<text>}| & \verb|\upshape|
			    & \upshape Upright shape \\
			意大利 & \verb|\textit{<text>}| & \verb|\itshape|
				& \itshape Italic shape \\
			倾斜& \verb|\textsl{<text>}| & \verb|\slshape|
				& \slshape Slanted shape \\
			小型大写 & \verb|\textsc{<text>}| & \verb|\scshape|
				& \scshape Small capitals shape \\
			\midrule
			字体系列 & 带参数的命令 & 申明命令 & 效果 \\
			中等 & \verb|\textmd{<text>}| & \verb|\mdseries|
				& \mdseries Medium series \\
			加宽加粗 & \verb|\textbf{<text>}| & \verb|\bfseries|
				& \bfseries Bold extended series \\
			\bottomrule
		\end{tabular}
	\end{table}

	中文字体没有英文字体那么复杂的成套的变体，
	各个字体之间一般都是独立的。
	因此，中文字体一般只用不同的字体族进行区分。
	ctex宏包和文档类（如ctexart）预定义了Windows常用的四种字体族：
	宋体、黑体、楷书、仿宋。
	为了方便使用，ctex宏包和文档类提供了简化命令：
	\verb|\songti| {\songti 宋体}，
	\verb|\heiti| {\heiti 黑体}，
	\verb|\kaishu| {\kaishu 楷书}，
	\verb|\fangsong| {\fangsong 仿宋}。
	
	\verb|\zihao{<字号>}|可以用来修改字体的大小，
	例如\verb|\zihao{4}|将字号设置为四号，
	字号前添加负号如\verb|\zihao{-4}|则表示小四号。
	
	除了修改字体和字号外，
	ctex宏包和文档类还增加了各种文字强调的方式，
	命令和效果如表 \ref{tab:fontemph} 所示。
	
	\begin{table}[hbp]
		\centering
		\caption{文字强调}
		\label{tab:fontemph}
		\begin{tabular}{ll}
			\toprule
			命令 & 效果\\
			\midrule
			\verb|\CJKunderdot{<文字>}| & \CJKunderdot{强调文字} \\
			\verb|\CJKunderline{<文字>}| & \CJKunderline{强调文字} \\
			\verb|\CJKunderdblline{<文字>}| & \CJKunderdblline{强调文字} \\
			\verb|\CJKunderwave{<文字>}| & \CJKunderwave{强调文字} \\
			\verb|\CJKsout{<文字>}| & \CJKsout{强调文字} \\
			\verb|\CJKxout{<文字>}| & \CJKxout{强调文字} \\
			\bottomrule
		\end{tabular}
	\end{table}
	
	\subsection{定义命令}
	\LaTeX 中的命令又称为宏，它们都以反斜线开头，
	一般格式为：
	
	\begin{quote}	% 引用环境：增加前后间距和缩进
		\begin{tabular}{ll}
			无参数： & \verb|\command| \\
			有n个参数： & \verb|\command{<arg1>}{<arg2>}...{<argn>}| \\
			有可选参数： & \verb|\command[<opt>]{<arg1>}{<arg2>}...{<argn>}| \\
		\end{tabular}
	\end{quote}

	例如，本文档在生成目录时就采用了命令\verb|\tableofcontents|。
	
	用户可以在导言区自定义命令，定义命令的基本格式为：
	
\begin{lstlisting}
\newcommand{<命令名>}[<参数个数>][<首参数默认值>]{<定义>}
\end{lstlisting}
	
	其中，命令名应当符合要求，以反斜线开头；
	当定义无参数命令时，
	参数个数和首参数默认值可连同中括号省略；
	当不给定默认值时，首参数默认值可连同中括号省略；
	当给定多个参数时，命令的定义中用井号加数字表示参数，
	如\verb|#2|表示第二个参数。
	
	例如我们可以定义一个命令来简化四元数乘法，
	按顺序输入两个四元数，
	输出四元数的乘法。
	定义为：
	
\begin{lstlisting}
\newcommand{\quadprod}[2]    % 可换行定义命令内容
    {\ensuremath{\mathfrak{#1} \otimes \mathfrak{#2}}}
\end{lstlisting}
	
	这样就定义了一个新命令\verb|\quadprod|。
	当使用\verb|\quadprod{q_b^a}{q_c^b}|时，
	编译后将输出\quadprod{q_b^a}{q_c^b}。
	
	除了用\verb|\newcommand|构造命令外，
	还可以使用\verb|\renewcommand|和\verb|\providecommand|。
	它们的用法相同。区别在于，
	\verb|\renewcommand|用于覆盖之前存在的同名命令；
	\verb|\providecommand|则会检测是否存在同名命令，
	若存在，则保留之前的用法，新定义不生效。
	
	\subsection{环境介绍}
	与命令相似，环境也可分为有参数的环境和无参数的环境。
	有参数环境的一般格式为：
	
\begin{lstlisting}
\begin{<环境名>}[<可选参数>]{<其他必要参数>} 
    <环境内容>
\end{<环境名>} 
\end{lstlisting}
	
	一个环境就是一个分组，它限定了一些命令的作用范围。
	除了使用环境，也可以用成对的花括号\{\}直接产生一个分组。
	
	我们给出一个分组的例子：
	为了将局部字体修改为黑体，
	只需要\verb|{\heiti 像这样}|构造分组，
	就能得到{\heiti 像这样}的局部黑体内容，
	而不会引起后面字体的改变。

	\section{编写文档}
	\subsection{页面设置}
	个人认为，页边距是使用 \LaTeX 编写文档时要考虑的第一个问题。
	因为如果最后修改页边距，
	一些所谓的盒子的内容有可能会超出边距，
	从而使文档需要重复调整。
	如果提前设置好边距，就能随时编译预览时发现问题。
	
	\LaTeX 对各种距离的定义较为复杂，
	为了能够快速修改，
	可以使用geometry宏包，
	并在导言区使用\verb|\geometry|命令对尺寸进行设置。
	
	如本文采用的设置为：
	
\begin{lstlisting}
\usepackage{geometry}
\geometry{ % 设置整体页面格式
	a4paper,
	left = 2.1cm,
	right = 2.1cm,
	bottom = 2.1cm,
	top = 2.5cm
}
\end{lstlisting}
	
	文档的第一页往往包含标题和作者等信息，
	也可以使用一个单独的封面。
	在标准文档类中，
	我们可以在正文中用\verb|\maketitle|生成标题页。
	该命令需要在导言区提前申明各种信息，如
	\verb|\title|、\verb|\author|等。
	
	除了使用预定义的标题页，
	我们也可以用titlepage环境自定义封面。
	该环境提供没有页码的单独一页，
	并使后面的内容页码从1开始计数。
	
	文档可使用\verb|\section|、\verb|\subsubsection|、
	\verb|\subsubsubsection|分别声明节、小节、小小节。
	\LaTeX 会自动进行编号。
	当需要使用到附录时，
	可使用\verb|\appendix|命令后，
	再分别使用\verb|\section|声明各级附录。
	附录格式的修改可使用appendix宏包，此处不具体说明。
	
	对于C\TeX 文档类
	\footnote{包括书籍ctexbook，报告ctexrep,文档ctexart。}，
	可以使用\verb|\CTEXsetup|命令来设置各章节标题的格式：
	
\begin{lstlisting}
\CTEXsetup[<选项1>=<值1>,<选项2>=<值2>...]{<对象类型>}
\end{lstlisting}

	对于ctexart文档，常用的对象类型包括section、subsection
	和subsubsection等。
	常用的选项如表 \ref{tab:ctexopt} 所示。
	
	\begin{table}
		\centering
		\caption{C\TeX 文档常用选项}
		\label{tab:ctexopt}
		\begin{tabular}{lll}
			\toprule
			选项 & 使用说明 & 值的示例 \\
			\midrule
			name & {<前名>,<后名>} & \verb|{第,节}| \\
			number & 设置编号的格式 & \verb|{\chinese{section}}| \\
			format & 章节名和章节标题的全局格式 & \verb|{\bfseries}| \\
			nameformat & 章节名和编号的格式 & 同format \\
			numberformat & 仅控制编号格式 & 同format \\
			titleformat & 仅控制章节标题格式 & 同format\\
			aftername & 章节名与标题之间的内容 & \verb|{\hspace{2ex}}| \\
			beforeskip & 章节标题前的段间距 & \verb|{\vspace{2em}}| \\
			afterskip & 章节标题后的段间距 & 同beforeskip \\
			indent & 章节标题的缩进长度 & 同aftername \\
			\bottomrule
		\end{tabular}
	\end{table}
	
	\subsection{添加目录}
	目录是最基本的自动化工具。
	\LaTeX 会自动收集章节命令所定义的各章节标题，
	用命令\verb|\tableofcontents|即可输出。
	类似地，命令\verb|\listoffigures|和\verb|\listoftables|
	会分别收集figure和table中\verb|\caption|命令的图表标题，
	产生图表的目录。
	为了将参考文献加入到目录中，
	可以使用tocbibind宏包。
	
	目录部分的页码常用大写的罗马数字表示，
	这可以使用命令\verb|\pagenumbering{<格式>}|进行修改，
	其中格式允许的值如表 \ref{tab:pagenum} 所示。
	
	在目录部分结束之后，
	常换页进入正文部分，
	并将页码计数重置为1，
	格式设置为阿拉伯数字，
	可采用以下命令：
	
\begin{lstlisting}
\clearpage					% 换页
\pagenumbering{arabic}		% 设置页码格式
\setcounter{page}{1}		% 正文重新开始于1
\end{lstlisting}
	
	\begin{table}[hbt]
		\centering
		\caption{页码数字格式}
		\label{tab:pagenum}
		\begin{tabular}{ll}
			\toprule
			格式 & 说明 \\
			\midrule
			arabic & 阿拉伯数字 \\
			roman & 小写的罗马数字 \\
			Roman & 大写的罗马数字 \\
			alph & 小写的字符形式 \\
			Alph & 大写的字符形式 \\
			\bottomrule
		\end{tabular}
	\end{table}
	
	\subsection{段落格式}
	在C\TeX 文档类中已经预设了正确的首行缩进，
	如果需要在某一段临时取消缩进，
	可以使用\verb|\noindent|命令；
	相反，为了在本来没有缩进的环境中临时设置缩进，
	可使用\verb|\indent|命令。
	这两个命令的作用范围均为一个段落，
	因而不需要再进行分组加以限定。
	
	除了段落的首行缩进，另一个关于分段的重要参数是段与段之间的垂直距离。
	这个距离由变量\verb|\parskip|控制，可以采用命令
	\verb|\setlength{\parskip}{<length>}|进行修改。
	这时应当注意，该命令会对后续的段落均产生影响。
	因此当只修改某些段落的段间距时，
	需要用花括号限制\verb|\setlength|的作用范围。
	
	在一些情况下，我们需要在文档内插入一段引文。
	quote环境常用来引用小段文字，
	该环境下没有首行缩进，
	且左右边距比正常文本稍微大一些，
	且增加了环境前后段的间距。
	比如下面给出一个使用这种环境的例子。
	
	\begin{quote}
		学而时习之，不亦说乎？
	\end{quote}
	
	而对于大段文字的引用，
	通常需要引入首行缩进，
	这时可采用quotation环境，
	如下面一段文字。
	
	\begin{quotation}
		小车正穿行在落基山脉蜿蜒曲折的盘山公路上。
		克里朵夫·李维静静地望着窗外，
		发现每当车子即将行驶到无路的关头，
		路边都会出现一块交通指示牌：
		“前方转弯！”或“注意！急转弯”。
		而拐过每一道弯之后，
		前方照例又是一片柳暗花明、豁然开朗。
		
		山路弯弯、峰回路转，
		“前方转弯”几个大字一次次地冲击着他的眼球，
		也渐渐叩开了他的心扉：
		原来，不是路已到了尽头，而是该转弯了。
		路在脚下，更在心中，心随路转，心路常宽。
		学会转弯也是人生的智慧，
		因为挫折往往是转折，危机同时是转机。
	\end{quotation}
	
	\subsection{定理环境}
	定理环境实际上是一类环境，在使用前需要在导言区进行定义：

\begin{lstlisting}
\newtheorem{thm}{定理}
\end{lstlisting}
	
	这样我们就得到了一个名为thm的定理类环境，
	它在使用时会自动产生形如“定理1”的提示。
	用同样的方法我们还可以定义引理、公理等。
	有时候我们希望定理的编号包含章节号，
	可在定义时增加参数即可：

\begin{lstlisting}
\newtheorem{thm}{定理}[section]
\end{lstlisting}
	
	定理环境还可以有一个可选参数，即定理的名字。
	下面给出这个环境的使用示例：
	
	\begin{thm}[勾股定理]
		直角三角形斜边的平方等于两直角边的平方和。
	\end{thm}	

	ntheorem宏包扩充了定理类环境的格式，
	在导言区使用\verb|\theoremstyle{<格式>}|可以方便地选择格式。
	可用的预定义格式如表 \ref{tab:thmstyle} 所示。
	
	\begin{table}
		\centering
		\caption{定理类环境格式}
		\label{tab:thmstyle}
		\begin{tabular}{ll}
			\toprule
			格式 & 说明 \\
			\midrule
			plain & 默认格式 \\
			break & 定理头换行 \\
			marginbreak & 编号在页边，定理头换行 \\
			changebreak & 定理头编号在前文字在后，换行 \\
			change & 定理头编号在前文字在后，不换行 \\
			margin & 编号在页边，定理头不换行 \\
			nonumberplain & 同plain格式，没有编号 \\
			nonumberbreak & 同break格式，没有编号 \\
			empty & 没有编号和定理名，只输出可选参数 \\
			\bottomrule
		\end{tabular}
	\end{table}
	
	该宏包具有很多设置命令和辅助功能，
	比如在使用该宏包时添加\verb|[thmmarks]|选项后，
	可以使用\verb|\theoremsymbol|命令在定理类环境末尾添加符号，
	这对定义证明环境表示证毕符号非常有用：
	
\begin{lstlisting}
% 导言区
\usepackage[thmmarks]{ntheorem}
{	% 利用分组使设置只对该分组内的定理类有效
	\theoremstyle{nonumberplain}
	\theoremheaderfont{\bfseries}
	\theorembodyfont{\normalfont}
	\theoremsymbol{\ensuremath{\Box}}
	\newtheorem{proof}{证明}
}
\end{lstlisting}
	
	这样，使用该定理类将会出现下面的效果：
	
	\begin{proof}
		当$x>0$时，将$\mathrm{e}^x$进行泰勒展开，有
		
		\[
			\mathrm{e}^x = \sum_{k=0}^{+\infty} \frac{x^k}{k!} 
			>\sum_{k=0}^{n} \frac{x^k}{k!}=P_n(x)
		\]
		
		证毕。
	\end{proof}
	
	\subsection{添加列表}
	\LaTeX 标准文档类提供了三种列表环境：
	带编号的enumerate环境、
	不编号的itemize环境
	和使用关键字的description环境。
	在列表环境内部使用\verb|\item|命令开始一个新的列表项，
	它可以带一个可选参数表示手动编号或关键字。
	
	这三种列表环境可以嵌套使用（最多四层），
	\LaTeX 会自动处理不同层次之间的缩进编号，
	例如：
	
	\begin{enumerate}
		\item 使用enumerate环境直接用\verb|\item|产生一个列表项
		\item[1b.] 使用\verb|\item[1b.]|测试手动编号
		\item 这里恢复正常的使用方法，并引入一次itemize嵌套
		\begin{itemize}
			\item 嵌套的第一层默认编号
			\item[\dag] 使用\verb|\item[\dag]|修改符号
			\item 恢复到默认设置
		\end{itemize}
	\end{enumerate}
	
	\subsection{使用颜色}
	可以利用xcolor宏包在文档中使用颜色。
	基本的命令包括声明式命令\verb|\color{<颜色>}|
	和带参数的命令\verb|\textcolor{<颜色>}{<文字>}|。
	前者对后面的内容都生效，
	而后者仅将参数部分的文字按指定的颜色输出。
	例如：
	
\begin{lstlisting}
{	% 使用分组，避免对后文上色
	\bfseries\color{red}		
	这是一段夹杂\textcolor{blue}{蓝色字符}的红色段落。
}
\end{lstlisting}
	
	{	% 使用分组，避免对后文上色
		\bfseries\color{red}		
		这是一段夹杂\textcolor{blue}{蓝色字符}的红色段落。
	}

	除了对文字上色，我们也能够对文字的“背景”进行上色。
	这用到所谓的“盒子”的功能，
	基本语法如下：
	
	\begin{quote}
		\verb|\pagecolor{<页面颜色>}| \\
		\verb|\colorbox{<盒子颜色>}{<文字>}| \\
		\verb|\fcolorbox{<线框颜色>}{<盒子颜色>}{<文字>}|
	\end{quote}
	
	为了直接使用颜色名如blue来表示颜色，
	可以在使用xcolor宏包时添加dvipsnames选项来扩充预定义的颜色名。
	除此之外，我们还有两种方式来定义新的颜色：
	
	\begin{quote}
		\verb|\definecolor{<色彩名>}{<模型>}{<分量值>}| \\
		\verb|\colorlet{<色彩名>}{<色彩标记>}|
	\end{quote}
	
	其中，色彩名是指新定义的颜色的名字。
	常用的颜色模型有：
	gray（灰度）、rgb（红绿蓝）和cmyk（印刷四分色）。
	xcolor提供了用已有颜色生成新颜色的色彩标记方法，
	常用的有：
	
	\begin{quote}
		\begin{tabular}{ll}
			半色调： & \verb|<颜色>!<百分数>| \\
			混合色： & \verb|<颜色>!<百分数>!<颜色>| \\
			互补色： & \verb|-<颜色>| \\
		\end{tabular}
	\end{quote}
	
	下面给出一些实例供参考：
	
\begin{lstlisting}
\begin{quote}
	\bfseries
	\textcolor{Blue!80}{80\%的蓝色} \\
	\textcolor{Red!60!Cyan}{60\%的红色加40\%的青色} \\
	\colorbox{Green}{\textcolor{-Green}{绿色背景与互补色文字}}
\end{quote}
\end{lstlisting}
	
	产生的效果为：
	
	\begin{quote}
		\bfseries
		\textcolor{Blue!80}{80\%的蓝色} \\
		\textcolor{Red!60!Cyan}{60\%的红色加40\%的青色} \\
		\colorbox{Green}{\textcolor{-Green}{绿色背景与互补色文字}}
	\end{quote}
	
	\subsection{使用图表}	
	\LaTeX 支持直接使用代码绘图，
	但这通常非常复杂。
	更常用的操作是利用其他软件绘图，
	然后将图片插入到文档之中。
	graphicx宏包提供了
	\verb|\includegraphics|命令来实现插图。
	通常情况下，
	我们会把插图单独放在figure环境中。
	例如，图 \ref{fig:railgun} 采用了以下代码：
	
\begin{lstlisting}
\begin{figure}[hbp]
	\centering
	\includegraphics[width=0.5\textwidth]{figure/mywife1.jpg}
	\caption{某春雨的超电磁炮}
	\label{fig:railgun}
\end{figure}
\end{lstlisting}
	
	\begin{figure}[hbp]
		\centering
		\includegraphics[width=0.7\textwidth]{figure/mywife1.jpg}
		\caption{某春雨的超电磁炮}
		\label{fig:railgun}
	\end{figure}
	
	从这个代码就可以看出，
	进入环境之后首先设置了在环境中生效的居中命令；
	通常情况下图片的名字在图片下面，
	所以先用\verb|\includegraphics|命令插入图片，
	并通过可选参数设置图片的宽度
	\footnote{常用选项有：宽度width、高度heigh和缩放比例scale}；
	再由\verb|\caption|生成标题；
	考虑到需要使用的交叉引用，
	接着用\verb|\label|打上标签。
	应当注意的是，
	为了产生正确的交叉引用，
	标签应当在\verb|\caption|命令之后。
	
	同样地，为了使用表格，
	可以使用在table环境中嵌套tabular环境。
	初次接触的小伙伴们可能会好奇，
	为什么制作表格还需要两层嵌套，
	而且这两个环境从含义上看来并没有很大的区别。
	我们可以这样认为：
	tabular环境用来制作表格，
	而table环境用来指明表格在文档中的位置。

	例如，表 \ref{tab:pagenum} 的代码如下：
	
\begin{lstlisting}
\begin{table}[hbp]
	\centering
	\caption{页码数字格式}
	\label{tab:pagenum}
	\begin{tabular}{ll}
		\toprule
		格式 & 说明 \\
		\midrule
		arabic & 阿拉伯数字 \\
		roman & 小写的罗马数字 \\
		Roman & 大写的罗马数字 \\
		alph & 小写的字符形式 \\
		Alph & 大写的字符形式 \\
		\bottomrule
	\end{tabular}
\end{table}
\end{lstlisting}

	首先我们来分析这个代码：
	我们在最外围申明了表格所处的table环境；
	然后设置该环境下内容居中显示；
	由于表格的标题一般在表格上方，
	所以先生成标题并打上标签；
	剩下的一大块便是由tabular环境构造的表格。
	
	然后我们讨论下tabular环境的使用方法：
	该环境有一个参数用来指明列格式，
	常用的列格式包括：
	
	\begin{itemize}
		\item \verb|l| --- 本列左对齐。
		\item \verb|c| --- 本列居中。
		\item \verb|r| --- 本列右对齐。
		\item \verb|p{<列宽>}| --- 指定列宽并允许自动换行。
		\item \verb=|= --- 画一条竖线，不占据表项计数。
		\item \verb|@{<内容>}| --- 任意添加内容，不占表项计数。
		\item \verb|*{<计数>}{<列格式说明>}| --- 将给定列格式按计数重复多次
	\end{itemize}

	在表内，采用符号\&指明对齐位置，
	并使用 \textbackslash\textbackslash 进行换行。
	表格中可以使用\verb|\hline|绘制横线，
	或使用\verb|\cline{<开始>-<结束>}|指定横线的范围。
	本文引用了booktabs宏包产生三线表的横线，
	可使用命令
	\verb|\toprule|、\verb|\midrule|、\verb|\bottomrule|
	分别绘制顶部横线、中间横线和底部横线。
	为了像\verb|\cline|那样绘制指定位置的横线，
	该宏包还提供了\verb|\cmidrule|命令。
	表格中绘制与单元格等高的竖线可用\verb|\vline|命令。
	
	在不单独指定列宽的情况下，
	tabular环境可以根据内容自动改变表格宽度，
	这在绝大多数情况下是非常好用的。
	然而有时候我们希望指定表格的总宽度，
	这就需要tabularx宏包提供的tabularx环境。
	该环境在指定列格式之前必须指定表格的总宽度。
	同时，该宏包提供了一个特殊的列格式X，
	它能够根据表的内容自动改变列宽，
	也可以和其他列格式一起使用。
	
	由于X格式默认是左对齐的，
	如果需要居中则需要\verb|\centering|进行设置。
	然而这个命令会影响单元格内的自动换行，
	需要加上命令\verb|\arraybackslash|进行恢复。
	如果这种格式比较常用，
	我们可以利用下面的命令自定义一个新的列样式。
	
\begin{lstlisting}
\newcolumntype{Y}{>{\centering\arraybackslash}X}
\end{lstlisting}
	
	有时表格需要用到单元格合并，
	对于列合并，可以直接使用下面的代码实现。

\begin{lstlisting}
\multicolumn{<列数>}{<列格式>}{<内容>}
\end{lstlisting}

	列合并的项数可以为1，
	常用于修改指定单元格的格式。
	表格的行合并相对复杂，
	需要利用multirow环境的同名命令，
	该命令的用法如下：
	
\begin{lstlisting}
\multirow{<行数>}{<列宽>}{<内容>}
\multirow{<行数>}*{<内容>}
\end{lstlisting}
	
	前一种格式会在内容达到指定列宽后自动换行，
	后一种格式会自动修改列宽，更为常用。
	
	如果在使用颜色宏包xcolor时添加table选项，
	就可以在表格中使用颜色。
	使用方法如表 \ref{tab:tabcolor} 所示。
	
	\begin{table}[hbp]
		\centering
		\caption{彩色表格的常用命令}
		\label{tab:tabcolor}
		\rowcolors{1}{SkyBlue!40}{Melon!40}
		\begin{tabular}{lll}
			命令 & 作用 & 常用位置 \\
			\verb|\columncolor{<颜色>}| & 改变列的颜色
				& \verb|>{<...>}|列格式说明符 \\
			\verb|\rowcolor{<颜色>}| &  改变行的颜色
				& 表格某一行的开头 \\
			\verb|\cellcolor{<颜色>}| & 改变单元格颜色
				& 单元格内 \\
			\verb|\rowcolors{<起始行>}{<奇数色>}{<偶数色>}|
				& 产生颜色交错的表格 & tabular等环境之前 \\
		\end{tabular}
	\end{table}
	
	以上介绍了插图和制表的基本命令，
	现在我们再来讨论下figure和table环境。
	细心的小伙伴一定发现上面两个例子在使用环境时
	均添加了\verb|[hbp]|的可选参数，
	这是什么意思呢？
	
	对于其他的环境，
	例如插入代码所使用的lstlisting环境，
	它在排版时保持了前后文相对关系不变，
	因而是固定的。
	而figure环境和table环境都是所谓的浮动体，
	它会根据当前页面的内容自动设置位置，
	从而避免页面出现大片的空白
	或图表超出页面范围。
	这两个环境的可选参数代表环境内容可以出现的位置，
	它们分别代表：
	
	\begin{itemize}
		\item \textbf{h} （here），浮动体可以放置在代码所在的上下文位置。
		\item \textbf{t} （top），浮动体可以放置在页面顶部，
			这可以是当前页面或下一页，
			当前页面排版时可能会出现在实际代码之前。
		\item \textbf{b} （bottom），浮动体可以放置在页面底部。
		\item \textbf{p} （page）一个或多个浮动体放置在单独的页面。
	\end{itemize}

	浮动体的允许位置可以由h、t、b、p任意组合，
	组合顺序不影响排版效果。
	figure和table的默认选项是\verb|[tbp]|。
	浮动体其他的一些预设值可能会限制浮动体内容出现的位置，
	可以在选项中使用！来临时取消预设值，
	如将可选参数设为\verb|[!htb]|。
	
	如果浮动体的自动排版达不到理想的效果，
	可以用float宏包对其进行设置。
	该宏包提供了一个新的位置选项H，
	它只能单独使用，
	其功能是取消环境的浮动。
	
	\subsection{数学公式}
	\AmS 相关宏包提供了很好的数学支持，
	这些宏包包括：amsmath、amsfonts、amssymb。
	针对数学公式的编辑在此不过多解释，
	有条件的小伙伴们可以在MathType中编辑，
	并设置其复制选项为AMSLaTeX代码。
	
	行内公式可以用美元符号表示，
	如\verb|$\frac{1}{2}$|的效果为$\frac{1}{2}$,
	\AmS 宏包会自动根据行内公式对字符进行调整。
	
	\AmS 宏包提供了很多环境用于在行间显示数学公式，
	这里仅对一些常用的环境进行介绍。
	其中，行间公式通常需要编号，
	且编号往往会带上章节号，
	这只需要在导言区添加\verb|\numberwithin{equation}{section}|即可。
	
	align环境提供了公式居中，编号靠右的默认格式，
	这也是一般科技文档的要求。
	当在该环境内编辑多行公式时，
	可使用\verb|&|指明对齐位置，
	并使用\verb|\\|进行换行。
	该环境默认对每一行公式进行编号，
	因此可以在每行公式后面为公式打上标签。
	如果不需要对公式进行编号，
	可使用\verb|\notag|取消编号。
	如果公式都不需要编号，
	可以使用align*环境。
	% 行间公式不需要空行来换行，否则会有多余空行
	\begin{align}
		{m_1}{{\ddot \vec r}_1} &=  
			-G\frac{{{m_1}{m_2}}}{{r_{12}^2}}\frac{{{{\vec r}_{21}}}}
			{{{r_{12}}}} 	\label{eq:r1}  \\
		{m_2}{{\ddot \vec r}_2} &=  
			- G\frac{{{m_1}{m_2}}}{{r_{12}^2}}\frac{{{{\vec r}_{12}}}}
			{{{r_{12}}}} 	\label{eq:r2}
	\end{align}
	
	观察这个示例公式，
	我们发现它们是同一类的，
	更好的表示方法是使它们具有相同的主编号，
	而增加子编号进行区分。
	这可以在align环境外嵌套subequations环境实现，如：
	
\begin{lstlisting}
\begin{subequations}
	\label{eq:rab}
	\begin{align}		% 公式代码由MathType生成
		{m_1}{{\ddot \vec r}_1} &=  
			-G\frac{{{m_1}{m_2}}}{{r_{12}^2}}
			\frac{{{{\vec r}_{21}}}}
			{{{r_{12}}}} 	\label{eq:ra}  \\
		{m_2}{{\ddot \vec r}_2} &=  
			- G\frac{{{m_1}{m_2}}}{{r_{12}^2}}
			\frac{{{{\vec r}_{12}}}}
			{{{r_{12}}}} 	\label{eq:rb}
	\end{align} 
\end{subequations}
\end{lstlisting}
	% continue
	\begin{subequations}
		\label{eq:rab}
		\begin{align}		% 公式代码由MathType生成
			{m_1}{{\ddot \vec r}_1} &=  
				-G\frac{{{m_1}{m_2}}}{{r_{12}^2}}
				\frac{{{{\vec r}_{21}}}}
				{{{r_{12}}}} 	\label{eq:ra}  \\
			{m_2}{{\ddot \vec r}_2} &=  
				- G\frac{{{m_1}{m_2}}}{{r_{12}^2}}
				\frac{{{{\vec r}_{12}}}}
				{{{r_{12}}}} 	\label{eq:rb}
		\end{align} 
	\end{subequations}

	这样，我们可以分别用\verb|\eqref{eq:rab}|
	引用主编号\eqref{eq:rab}，
	也可以用\verb|\eqref{eq:ra}|引用子编号\eqref{eq:ra}。
	
	有时我们还会用多行公式展开推导的详细过程，
	这时我们不仅希望公式按等号进行对齐，
	还希望整个推导只占有一个编号。
	这时可以使用split环境，
	该环境的用法与align相同。例如：
	
\begin{lstlisting}
\begin{align}
	\label{eq:rc}
	\begin{split}
		{{\ddot \vec r}_C} &= \frac{{{m_1}}}{{{m_1} + {m_2}}}
			{{\ddot \vec r}_1} + \frac{{{m_2}}}{{{m_1} + {m_2}}}
			{{\ddot \vec r}_2} \\ 
		&= \frac{1}{{{m_1} + {m_2}}}\left[ {\left( { - G
			\frac{{{m_1}{m_2}}}{{r_{12}^2}}
			\frac{{{{\vec r}_{21}}}}{{{r_{12}}}}} \right) + 
			\left( { - G\frac{{{m_1}{m_2}}}{{r_{12}^2}}
			\frac{{{{\vec r}_{12}}}}{{{r_{12}}}}} 
			\right)} \right] \\ 
		&= \frac{{ - G{m_1}{m_2}}}{{\left( {{m_1} + {m_2}}
			\right)r_{12}^3}}\left( {{{\vec r}_{21}} + 
			{{\vec r}_{12}}} \right) \\ 
		&= \vec 0 
	\end{split}
\end{align}
\end{lstlisting}
	% continue
	\begin{align}
		\label{eq:rc}
		\begin{split}
			{{\ddot \vec r}_C} &= \frac{{{m_1}}}{{{m_1} + {m_2}}}
			{{\ddot \vec r}_1} + \frac{{{m_2}}}{{{m_1} + {m_2}}}
			{{\ddot \vec r}_2} \\ 
			&= \frac{1}{{{m_1} + {m_2}}}\left[ {\left( { - G
				\frac{{{m_1}{m_2}}}{{r_{12}^2}}
				\frac{{{{\vec r}_{21}}}}{{{r_{12}}}}} \right) + 
				\left( { - G\frac{{{m_1}{m_2}}}{{r_{12}^2}}
				\frac{{{{\vec r}_{12}}}}{{{r_{12}}}}} 
				\right)} \right] \\ 
			&= \frac{{ - G{m_1}{m_2}}}{{\left( {{m_1} + {m_2}}
				\right)r_{12}^3}}\left( {{{\vec r}_{21}} + 
				{{\vec r}_{12}}} \right) \\ 
			&= \vec 0 
		\end{split}
	\end{align}
	
	\subsection{插入代码}
	\LaTeX 输入特殊符号时不太便利，
	然而有时候我们必须经常性地使用特殊符号，
	例如在排版计算机程序源代码的时候。
	此时需要使用抄录功能。
	
	使用命令\verb|\verb<符号><抄录内容><符号>|
	可以将抄录内容原封不动地输出到正文中，
	两端的符号可以任意给定，但是要确保相同。
	例如\verb~\verb|\TeX|~采用“|”作为边界。
	
	在插入大段代码时，可以使用listings宏包提供的lstlisting环境。
	但是直接使用往往得不到很好的效果，
	通常我们需要在导言区使用\verb|\lstset|进行设置。
	例如本文的设置如下：

\begin{lstlisting}
% 导言区
\usepackage{listings}
\usepackage{listings}
\lstset{ % 代码环境整体设置
	basicstyle = \ttfamily,
	keywordstyle = \bfseries,
	commentstyle = \rmfamily\upshape,
	stringstyle = \ttfamily\slshape,
	tabsize = 4,
	backgroundcolor = \color{lightgray},	% 需先引入颜色宏包
	frame = single,
	language = TeX							% 设置默认语言
}
\end{lstlisting}
	
	需要说明的是，lstlisting环境对缩进是敏感的。
	这就是说，从\verb|\begin{lstlisting}|开始的缩进都将进入正文。
	该环境还可以使用可选参数，
	对环境进行临时修改。
	为了使代码根据不同语言实现高亮，
	可以在使用该环境使添加可选参数language进行设置，
	例如上面的例子就引入了参数\verb|[language=TeX]|。
	
	\subsection{插入引用}
	最简单的引用便是使用\verb|\footnote{<脚注内容>}|产生脚注，
	例如这个位置\footnote{这是一个脚注的例子}。
	脚注是自动编号的，也可以使用可选参数修改编号，
	但是这不改变原来脚注的编号。
	
	在例如表格等环境中，\verb|\footnote|命令往往不能直接使用。
	这时我们可以把脚注标记和脚注内容分开\footnotemark，
	这两条命令分别为
	\verb|\footnotemark|和\verb|\footnotetext{<脚注内容>}|。
	\footnotetext{这又是一个脚注}
	
	\LaTeX 默认的脚注格式不太好看，
	我们可以在导言区进行以下设置进行优化。
	
\begin{lstlisting}
\usepackage[perpage]{footmisc}		% 脚注每页清零
\usepackage{pifont}					% 优化带圈的数字
\renewcommand{\thefootnote}{\ding{\numexpr171+\value{footnote}}}
\end{lstlisting}
	
	交叉引用可以通过一个符号标签引用文档中某个对象的
	编号、页码、或标题等信息，
	而不必知道这个对象具体在什么位置。
	我们需要两个步骤来实现交叉引用：
	定义标签和引用标签。
	
	定义标签是在合适的位置给一个带参数的对象添加标签，
	命令为\verb|\label{<标签>}|。
	标签的名字最好是简洁而有用的名字，
	不能包含特殊字符。
	标签的位置常在\verb|\section|命令之后；
	或table和figure环境中\verb|\caption|命令之后；
	也可以在公式之后。
	
	创建标签后，可分别使用
	\verb|\ref|、\verb|\pageref|、\verb|\nameref|
	引用标签对象的计数、所在页码和名称。
	对于公式的引用，通常会给编码加上英文括号，
	好在 \AmS 的宏包在支持数学的基础上
	增加了命令\verb|\eqref|专门用于公式引用，
	所以公式的引用会略有不同。
	
	例如，我们使用下面的代码进行交叉引用的测试：

\begin{lstlisting}
% 前面已使用 \label{tab:ctexopt} 定义标签
表 \ref{tab:ctexopt} 在第 \pageref{tab:ctexopt} 页，
名字是 \nameref{tab:ctexopt}。
\end{lstlisting}
	
	输出为：表 \ref{tab:ctexopt} 在第 \pageref{tab:ctexopt} 页，
	名字是 \nameref{tab:ctexopt}。
	
	细心的朋友对比后会发现，
	\verb|\tableofcontents|生成的目录不是超级链接，
	而且编译生成的pdf文档中不会出现标签。
	为了解决这个问题，
	只需要使用hyperref宏包。
	或者在申明文档类时加上可选项，
	如下
	
\begin{lstlisting}
% 导言区
\documentclass[hyperref,UTF8]{ctexart}
\end{lstlisting}
	
	同时，该宏包还提供了命令
	\verb|\url{<URL>}|和\verb|\href{<URL>}{<文字>}|
	产生超级链接。
	比如，这里我可以班门弄斧一下
	\href{https://ichunyu.github.io}{我的博客}。
	
	除了上面这些引用，
	科技类文档中最重要的引用就是文献引用了。
	这类引用首先需要一个数据库文件，
	其通常以bib为后缀。
	该文件由多项参考文献构成，
	这里给出其中的一个例子：
	
\begin{lstlisting}
@article{canuto_embedded_2018,
title = {Embedded model control: {Reconciling} 
	modern control theory and error-based control design},
volume = {16},
issn = {2198-0942},
shorttitle = {Embedded model control},
url = {https://doi.org/10.1007/s11768-018-8130-1},
doi = {10.1007/s11768-018-8130-1},
language = {en},
number = {4},
urldate = {2019-07-16},
journal = {Control Theory and Technology},
author = {Canuto, Enrico and Novara, Carlo and Colangelo, Luigi},
month = nov,
year = {2018},
keywords = {embedded model control, disturbance rejection,
	error loop, error-based design, Modern control theory},
pages = {261--283}
}
\end{lstlisting}
	
	通常情况下bib文件不需要我们手动编写，
	文献检索网站和各种文献管理软件一般都支持这种BibTeX格式的导出。
	
	文献的引用需要以下三个步骤：
	
	\begin{enumerate}
		\item 使用\verb|\bibliographystyle|命令设定参考文献的格式，
		这通常在导言区完成。基本的文献格式有：
		plain --- 按作者、日期、标题排序；
		unsrt --- 不排序但保持引用次序；
		alpha --- 使用一种三字母缩写的方式编号并按作者排序；
		abbrv --- 与plain基本相同，但是定义了一些缩写。
		
		\item 在正文中使用\verb|\cite|命令引用所需要的文献，
		当引用多个文献时用逗号隔开。
		
		\item 使用\verb|\bibliography|命令指明要使用的文献数据库，
		即包含文献信息的bib文件。
		同时， \LaTeX 会在这个命令的位置插入参考文献列表。
	\end{enumerate}

	例如，使用以下命令进行文献引用：

\begin{lstlisting}
Canuto教授提出了一种模型嵌入控制\cite{canuto_embedded_2018}。
\end{lstlisting}
	
	其效果为：Canuto教授提出了一种模型嵌入控制\cite{canuto_embedded_2018}。
	参考文献列表见本文末尾。
	
	默认的文献样式不能满足所有出版社的要求，
	这时我们可以使用natbib宏包进行修改。
	该宏包同时定义三种专用的格式：
	plainnat、abbrvnat、unstrnat。
	
	natbib宏包提供了\verb|\setcitestyle|命令来设置引用命令的输出格式，
	在参数中可以设置以下选项：
	
	\begin{itemize}
		\item 选择引用模式：
		authoryear代表作者年代模式，如[Canuto,2018];
		numbers表示数字序号模式，如[3];
		super表示数字上标模式，如${}^{[64]}$。
		
		\item 括号：
		round圆括号，
		square方括号，
		或是用\verb|open={<左括号>}|和\verb|close={<右括号>}|
		分别进行设置。
		
		\item 多个引用之间的标点：
		semicolon分号，comma逗号，
		或是用\verb|citesep={<符号>}|进行设置。
		
		\item 作者与年代之间的符号：\verb|aysep={<符号>}|。
		
		\item 同一作者的几个年代间的符号\verb|yysep={<符号>}|。
		
		\item 在引用命令可选参数的说明文字前的符号\verb|notesep={<符号>}|。
	\end{itemize}

	简单起见，我们可以在使用宏包时添加选项来快速修改格式，如下：
	
\begin{lstlisting}
\usepackage[super,square]{natbib}
\end{lstlisting}
		
	\section{使用模板}
	除了利用各种宏包修改格式编写自己的文档外，
	以期刊文章为主的各类科技文档都可以找到相应的模板，
	这些模板针对要求对格式进行了修改，
	获得模板后基本可以直接使用。
	
	需要注意的使，
	一些模板为了使用上的方便，
	重新定义了一些命令，
	其使用方法可能略有差别。
	不过这一般会在模板的使用说明中指明，
	小伙伴们在使用模板前应当注意阅读使用说明。
	
	最后，祝大家学习愉快！
	
	Happy \LaTeX ing！
	
	\bibliography{ref}
	
%	\appendix
%	\section{本文使用的宏包}
%	\begin{table}
%		\centering
%		\caption{本文使用的宏包汇总}
%		\label{tab:pkgs}
%		\begin{tabular}{ll}
%			\toprule
%			宏包 & 说明 \\
%			\midrule
%			geometry & 页边距相关设置 \\
%			hyperref & 使用超链接功能 \\
%			tocbibind & 添加参考文献到目录 \\
%			（\AmS 宏包） & 数学相关宏包：amsmath，amsfonts， amssymb  \\
%			metalogo & \XeLaTeX 等logo \\
%			xcolor & 颜色相关宏包 \\
%			listings & 支持语法高亮的代码环境 \\
%			footmisc & 脚注编号每页清零 \\
%			pifont & 优化脚注带圆圈的数字 \\
%			graphicx & 插图功能 \\
%			tabularx & 定宽表格 \\
%			multirow & 表格的行合并 \\
%			booktabs & 三线表线型 \\
%			float & 浮动体控制 \\
%			natbib & 参考文献格式 \\
%			ntheorem & 定理类格式 \\
%			\bottomrule
%		\end{tabular}
%	\end{table}

\end{document}
